\chapter{Layer agentico e MCP}
\label{chap:agentico}

\section{Ruoli dell'agente}
\label{sec:ruoli_agente}

Il sistema realizza i seguenti ruoli operativi distinti, in linea con il requisito
di almeno due responsabilit\`a agentiche separate:

\begin{enumerate}[leftmargin=*]
    \item \textbf{Content interpretation}: estrazione semantica dal corpo dell'email e
          dagli allegati (PDF, DOCX, XLSX, testo semplice) tramite il parser
          \texttt{eml\_parser} e Apache Tika.

    \item \textbf{Document classification \& filing}: associazione di ogni allegato alla
          cartella pi\`u appropriata dell'albero del tenant, accompagnata da
          \texttt{confidence} (0.0--1.0) e \texttt{reasoning} testuale. Il risultato
          pu\`o essere un'azione autonoma (\texttt{auto\_filed}) o una proposta
          (\texttt{in\_inbox}).

    \item \textbf{Operational summarisation}: produzione di descrizioni human-readable
          e \emph{details} strutturati (JSONB) per ogni operazione, che alimentano
          la timeline operativa visibile a utente e auditor.

    \item \textbf{Escalation policy}: instradamento automatico all'\emph{inbox}
          quando la confidenza non supera la soglia, con creazione di una proposta
          reversibile per la revisione umana.

    \item \TODO{Monitoring operativo del cluster Kubernetes tramite un secondo agente
          connesso via MCP a metriche Prometheus e log Loki --- da implementare}
\end{enumerate}

\section{Perch\'e un agente e non un classificatore deterministico}
\label{sec:perche_agente}

La scelta di un componente agentico sopra un classificatore tradizionale
(\emph{rule-based} o ML su feature ingegnerizzate) \`e motivata da tre osservazioni:

\begin{itemize}[leftmargin=*]
    \item \textbf{L'albero delle cartelle \`e idiosincratico al tenant e cambia nel tempo}:
          un classificatore con etichette fisse richiederebbe ri-training continuo per
          ogni cliente, con costi operativi insostenibili.
    \item \textbf{Le evidenze disponibili sono eterogenee}: soggetto e corpo dell'email,
          MIME type, nome file, anteprima testuale del documento. Gli LLM eccellono
          nel ragionamento \emph{few-shot} su input miscellanei.
    \item \textbf{Evolvibilit\`a zero-shot}: l'aggiunta di nuove cartelle non richiede
          modifiche al codice. L'agente le scopre autonomamente via MCP alla successiva
          invocazione.
\end{itemize}

La componente deterministica \`e comunque presente e irrinunciabile: estrazione
attachment, hashing SHA-256, threshold di confidenza, validazione UUID, scrittura DB.
Queste operazioni separano con nettezza il piano probabilistico dal piano deterministico.

\section{Parsing degli allegati (\texttt{eml\_parser})}
\label{sec:eml_parser}

Il modulo \texttt{eml\_parser} estrae anteprime testuali controllate per ciascun MIME type:

\begin{itemize}[leftmargin=*]
    \item \textbf{PDF}: prime e ultime 2~pagine, cap di 1000~caratteri per lato con
          separatore \texttt{[...]}. Ritorna \texttt{None} per PDF scansionati senza
          testo estraibile (condizione che segnala all'agente di abbassare la confidenza).
    \item \textbf{DOCX}: testa e coda di 1000~caratteri sul testo concatenato dei paragrafi.
    \item \textbf{XLSX}: prime 30~righe di ogni foglio (max 3~fogli) come TSV, capate
          a 2000~caratteri.
    \item \textbf{Plain text}: primi 500~caratteri.
\end{itemize}

Questa logica \`e intenzionalmente \emph{conservativa}: l'agente vede solo un sottoinsieme
deterministico, riproducibile e di dimensione bounded. Tale disciplina costituisce il
primo livello di \emph{economic guardrail} (limite ai token in input).

\section{Atomic Agents e structured output}
\label{sec:atomic_agents}

Il filing agent \`e implementato con il framework \textbf{Atomic Agents} e la libreria
\textbf{instructor} come adattatore per forzare l'output del modello LLM a conformarsi
allo schema Pydantic. Gli schema input/output sono:

\begin{lstlisting}[language=Python, caption={Schema input e output del filing agent (\texttt{agent\_worker/agent.py}).}]
class AttachmentInfo(BaseIOSchema):
    filename: str
    mime_type: str
    size_bytes: int
    text_preview: str | None

class AttachmentDecision(BaseModel):
    filename: str
    folder_id: str | None  # UUID dalla folder tree, o null
    confidence: float      # ge=0.0, le=1.0
    reasoning: str

class FilingInput(BaseIOSchema):
    context_title: str           # soggetto email o nome file
    context_source: str          # mittente o "(manual upload)"
    context_description: str | None  # corpo email (max ~2000 char)
    attachments: list[AttachmentInfo]

class FilingOutput(BaseIOSchema):
    decisions: list[AttachmentDecision]  # una per allegato
\end{lstlisting}

La constrained generation fornita da \texttt{instructor} trasforma l'output del modello
in un \texttt{FilingOutput} validato. Risposte strutturalmente non conformi vengono
scartate, eliminando la classe di errori da \emph{free-form text}.

\subsection{System prompt}

Il system prompt \`e strutturato in \texttt{background}, \texttt{steps} e
\texttt{output\_instructions}, con la folder tree iniettata come
\texttt{FolderTreeContextProvider}:

\begin{lstlisting}[language=Python, caption={System prompt del filing agent (estratto da \texttt{agent.py}).}]
SystemPromptGenerator(
  background=[
    "You are an expert document archivist for a DMS.",
    "Your task is to classify email attachments and assign each one "
    "to the most appropriate folder.",
    "The available folder structure is provided in the context below.",
  ],
  steps=[
    "Read the email subject, sender, and body to understand context.",
    "Pay special attention to the email body: it often describes "
    "what the attachments are about.",
    "For each attachment, analyze filename, MIME type, size, and "
    "any text preview.",
    "If no folder is a good match, set folder_id to null.",
    "Set confidence to 1.0 only if certain. Lower it for ambiguous cases.",
  ],
  output_instructions=[
    "Return exactly one decision per attachment in 'decisions'.",
    "folder_id must be a valid UUID from the folder tree, or null.",
    "confidence must be a float between 0.0 and 1.0.",
  ],
  context_providers={"folders": FolderTreeContextProvider(folder_tree)},
)
\end{lstlisting}

\section{Design del MCP Server}
\label{sec:mcp_design}

Il MCP server espone tre strumenti, raggruppati in due \emph{capability groups}:

\begin{enumerate}[leftmargin=*]
    \item \textbf{Folder structure inspection} (gruppo 1):
          \texttt{list\_folders(user\_id, parent\_id?)},
          \texttt{get\_folder(user\_id, folder\_id)}.
    \item \textbf{Document inventory inspection} (gruppo 2):
          \texttt{list\_documents(user\_id, folder\_id?)}.
\end{enumerate}

\begin{lstlisting}[language=Python, caption={Tool \texttt{list\_folders} del MCP server (\texttt{mcp\_server/server.py}).}]
@mcp.tool()
def list_folders(user_id: str, parent_id: str | None = None) -> list[dict]:
    """List folders owned by user_id, root or sub-folders."""
    db = SessionLocal()
    try:
        query = db.query(Folder).filter(Folder.owner_id == UUID(user_id))
        if parent_id:
            query = query.filter(Folder.parent_id == UUID(parent_id))
        else:
            query = query.filter(Folder.parent_id.is_(None))
        folders = query.order_by(Folder.name).all()
        return [{"id": str(f.id), "name": f.name,
                 "parent_id": str(f.parent_id) if f.parent_id else None}
                for f in folders]
    finally:
        db.close()
\end{lstlisting}

Tutti gli strumenti sono \textbf{strettamente read-only}. Non esiste nessun tool MCP
di scrittura, eliminazione o spostamento. Ogni tool call \`e \emph{tenant-scoped}
tramite il filtro \texttt{owner\_id}, impedendo che una prompt-injection induca
l'agente a leggere dati di un altro tenant.

Il server usa \textbf{FastMCP} con trasporto \textbf{Server-Sent Events} (SSE),
esposto sulla porta 8001 e raggiungibile solo dall'agent worker sulla rete Docker interna.

\section{Pre-fetch della folder tree}
\label{sec:prefetch}

L'agent worker pre-fetcha l'intera folder tree con \texttt{MCPClient.build\_folder\_tree}
e la inietta nel system prompt come contesto statico prima di invocare l'LLM. La scelta
del \emph{pre-fetch} (piuttosto che tool call iterativi durante l'inferenza) \`e motivata
da:

\begin{itemize}[leftmargin=*]
    \item \textbf{Determinismo}: l'agente vede una snapshot consistente dell'albero,
          evitando race con creazioni di cartelle concorrenti.
    \item \textbf{Prevedibilit\`a del costo}: il numero di token in input dipende solo
          dalla dimensione dell'albero del tenant, condizione necessaria per imporre
          un cost ceiling.
    \item \textbf{Riduzione delle round-trip}: una sola materializzazione vs.\
          $O(N)$ tool call durante l'inferenza.
\end{itemize}

\section{Pipeline di classificazione end-to-end}
\label{sec:pipeline}

Il modulo \texttt{agent\_worker/worker.py} implementa due pipeline parallele:

\paragraph{Pipeline email.}
Richiamata da \texttt{\_process\_job}: scarica \texttt{.eml} $\rightarrow$ parsa
email ($\rightarrow$ estrae allegati $\rightarrow$ carica su S3 $\rightarrow$ crea
\texttt{AgentFile}) $\rightarrow$ indica in OpenSearch $\rightarrow$ invoca LLM
$\rightarrow$ decide auto\_filed o in\_inbox.

\paragraph{Pipeline manual upload.}
Richiamata da \texttt{\_process\_manual\_upload}: scarica il file da S3 $\rightarrow$
Tika estrae il testo $\rightarrow$ indicizza in OpenSearch $\rightarrow$ se
\texttt{needs\_classification = True}: invoca LLM $\rightarrow$ decide; altrimenti
imposta status \texttt{indexed}.

\begin{lstlisting}[language=Python, caption={Decisione auto-filing nel worker (\texttt{agent\_worker/worker.py}).}]
folder_uuid = UUID(decision.folder_id) if decision.folder_id else None

if folder_uuid and decision.confidence >= settings.auto_file_threshold:
    agent_file.status = "auto_filed"
    agent_file.auto_filed = True
    doc.folder_id = folder_uuid
    _log_op(db, user_id=user.id, op_type="auto_filed",
            description=f"Archiviato automaticamente '{att.filename}' "
                        f"(confidence: {decision.confidence:.0%})",
            details={"folder_id": str(folder_uuid),
                     "confidence": decision.confidence})
else:
    agent_file.status = "in_inbox"
    _log_op(db, user_id=user.id, op_type="sent_to_inbox",
            description=f"'{att.filename}' inviato in inbox per revisione "
                        f"(confidence: {decision.confidence:.0%})")
\end{lstlisting}

\section{Indicizzazione in OpenSearch}
\label{sec:opensearch_indexing}

Ogni documento elaborato viene indicizzato in OpenSearch con il seguente mapping:

\begin{lstlisting}[language=Python, caption={Mapping OpenSearch (\texttt{agent\_worker/worker.py}).}]
OPENSEARCH_INDEX_MAPPING = {
    "settings": {"index": {"knn": True}},
    "mappings": {"properties": {
        "document_id": {"type": "keyword"},
        "owner_id":    {"type": "keyword"},
        "folder_id":   {"type": "keyword"},
        "filename":    {"type": "text",
                        "fields": {"keyword": {"type": "keyword"}}},
        "mime_type":   {"type": "keyword"},
        "text":        {"type": "text", "analyzer": "standard"},
        "embedding":   {"type": "knn_vector",
                        "dimension": 384,
                        "method": {"name": "hnsw",
                                   "space_type": "cosinesimil",
                                   "engine": "lucene"}},
        "created_at":  {"type": "date"},
    }}
}
\end{lstlisting}

L'indice \`e \emph{per-tenant}: ogni utente ha il proprio indice
\texttt{documents\_\{user\_id\_senza\_trattini\}}. Un cache in memoria
(\texttt{\_existing\_indices: set[str]}) evita chiamate HEAD ripetute per verificarne
l'esistenza.

\section{Secondo agente: monitoring Kubernetes}
\label{sec:secondo_agente}

\TODO{Questa sezione descriver\`a il secondo agente agentico, responsabile del monitoring
operativo del cluster Kubernetes. L'agente si connetterà via MCP a un secondo MCP server
che espone tool read-only su metriche Prometheus (CPU, memoria, latenza per servizio) e
log Loki. L'agente interpreterà anomalie, proporrà azioni di remediation e le sottoporrà
all'operatore per approvazione. Nessuna azione distruttiva sarà automatizzata. Ogni
proposta sarà loggata nell'audit trail.}
